\documentclass[11pt]{article}
\usepackage[margin=1in]{geometry}
\usepackage[table]{xcolor}
\usepackage{booktabs}
\usepackage{siunitx}
\usepackage{multirow}
\usepackage{tabularx}

% Custom brand color defined in the preamble (tests preamble reuse)
\definecolor{tx2036blue}{RGB}{0,84,135}
\definecolor{stripegray}{RGB}{235,240,245}

% Custom macro used inside table cells (tests preamble reuse)
\newcommand{\hi}[1]{\textcolor{tx2036blue}{\textbf{#1}}}

\sisetup{
  group-separator = {,},
  group-minimum-digits = 4,
  table-align-text-after = false
}

\title{Regional Workforce and Housing Pressures in Texas:\\A Styled Tables Demonstration}
\author{Center for Applied Policy Illustration}
\date{July 2026}

\begin{document}
\maketitle

\section{Overview}

Texas continues to add residents faster than nearly every other state, and
the composition of that growth varies sharply across its metropolitan and
rural regions. This short report illustrates how regional labor markets,
housing costs, and public investment interact, using illustrative figures
assembled for demonstration purposes only. None of the numbers in the
tables that follow should be cited as official statistics; they exist to
exercise a table-rendering pipeline.

Section~\ref{sec:labor} reviews labor-market indicators, drawing on the
striped summary in Table~\ref{tab:labor}. Section~\ref{sec:precision}
turns to precisely aligned fiscal measures in
Table~\ref{tab:fiscal}. Section~\ref{sec:grid} presents a cross-tabulated
grid of program participation in Table~\ref{tab:grid}, and
Section~\ref{sec:housing} closes with a full-width narrative comparison
of housing conditions in Table~\ref{tab:housing}.

\section{Labor-market indicators}\label{sec:labor}

Employment growth in the state's mid-sized metros has outpaced the large
urban triangles in percentage terms, even as absolute job gains remain
concentrated in the biggest markets. Table~\ref{tab:labor} summarizes
five illustrative regions with alternating row shading to aid scanning.
Readers should note that the strongest percentage gains appear in the
smaller eastern markets, a pattern consistent with lower baseline
employment counts rather than unusually strong demand.

\begin{table}[htbp]
  \centering
  \caption{Illustrative regional employment indicators, striped}
  \label{tab:labor}
  \rowcolors{2}{stripegray}{white}
  \begin{tabular}{lrrr}
    \toprule
    Region & Jobs added & Growth (\%) & Openings ratio \\
    \midrule
    Nacogdoches Basin   & 8{,}214  & \hi{6.7} & 1.93 \\
    Palo Duro Corridor  & 3{,}471  & 2.9      & 1.18 \\
    Balcones Rim        & 21{,}057 & 4.1      & 1.42 \\
    Caprock Flats       & 1{,}663  & \hi{7.3} & 2.07 \\
    Neches Delta        & 5{,}829  & 3.6      & 1.31 \\
    \bottomrule
  \end{tabular}
\end{table}

\section{Fiscal measures with decimal alignment}\label{sec:precision}

Comparing per-capita spending across regions requires careful decimal
alignment, particularly when values span more than one order of
magnitude. Table~\ref{tab:fiscal} uses fixed-format numeric columns so
that magnitudes can be compared at a glance. The variation in
transportation outlays is the widest of the three categories shown,
reflecting the long highway mileage maintained in sparsely populated
areas.

\begin{table}[htbp]
  \centering
  \caption{Illustrative per-capita outlays by category (dollars)}
  \label{tab:fiscal}
  \begin{tabular}{l S[table-format=4.2] S[table-format=3.2] S[table-format=4.1]}
    \toprule
    Region & {Education} & {Health} & {Transport} \\
    \midrule
    Nacogdoches Basin  & 4187.53 & 612.44 & 1093.7 \\
    Palo Duro Corridor & 3958.21 &  587.09 & 2841.6 \\
    Balcones Rim       & 4402.86 & 634.71 &  876.2 \\
    Caprock Flats      & 3719.44 & 559.83 & 3164.9 \\
    Neches Delta       & 4066.02 & 601.57 & 1428.3 \\
    \bottomrule
  \end{tabular}
\end{table}

\section{Program participation grid}\label{sec:grid}

Participation in workforce training programs differs by both region and
age band. Table~\ref{tab:grid} cross-tabulates enrollment counts,
grouping the two younger cohorts and the two older cohorts under spanning
headers, with region groups collapsed into two multirow blocks. The
younger cohorts dominate enrollment in every region, though the margin
narrows in the rural block.

\begin{table}[htbp]
  \centering
  \caption{Illustrative training enrollment by region group and age band}
  \label{tab:grid}
  \begin{tabular}{llrrrr}
    \toprule
    & & \multicolumn{2}{c}{Younger cohorts} & \multicolumn{2}{c}{Older cohorts} \\
    \cmidrule(lr){3-4} \cmidrule(lr){5-6}
    Group & Region & {18--24} & {25--34} & {35--49} & {50--64} \\
    \midrule
    \multirow{2}{*}{Metro} & Balcones Rim  & 4{,}912 & 6{,}358 & 3{,}147 & 1{,}026 \\
                           & Neches Delta  & 1{,}744 & 2{,}209 & 1{,}183 &   417 \\
    \midrule
    \multirow{2}{*}{Rural} & Caprock Flats &   386   &   541   &   472   &   298 \\
                           & Palo Duro     &   659   &   803   &   744   &   511 \\
    \bottomrule
  \end{tabular}
\end{table}

\section{Housing conditions}\label{sec:housing}

Housing affordability now ranks among the top concerns reported by
employers recruiting to the state's fast-growing regions. Because the
qualitative assessments are longer than a numeric cell can comfortably
hold, Table~\ref{tab:housing} spreads flexible-width columns across the
full text width. The pattern is familiar: where job growth is fastest,
permitting has not kept pace, and median rents have climbed accordingly.

\begin{table}[htbp]
  \centering
  \caption{Illustrative housing conditions, full-width narrative comparison}
  \label{tab:housing}
  \begin{tabularx}{\textwidth}{lXXr}
    \toprule
    Region & Supply assessment & Policy response & Median rent \\
    \midrule
    Nacogdoches Basin & Permitting lags demand by roughly eleven months; infill parcels scarce near the mill district & Fee waivers piloted for accessory dwellings in \hi{2024} & \$1{,}287 \\
    Balcones Rim & Multifamily pipeline robust downtown but thin in the western foothill tracts & Density bonus adopted along the Ridgeline transit spine & \$1{,}942 \\
    Caprock Flats & Aging stock dominates; vacancy concentrated in pre-1970 units needing rehabilitation & Rehab grant fund capitalized at \$4.6 million & \$743 \\
    \bottomrule
  \end{tabularx}
\end{table}

\section{Conclusion}

The four tables above (Tables~\ref{tab:labor}--\ref{tab:housing})
illustrate a common analytic arc: scan-friendly summaries, precisely
aligned fiscal detail, cross-tabulated program data, and narrative
comparisons. A production report would pair each with sourcing notes and
uncertainty ranges; here the figures are intentionally synthetic.

\end{document}
